% Compile twice: lualatex 30003377.tex
\documentclass[11pt]{article}
\usepackage[margin=25mm]{geometry}
\usepackage{fontspec,unicode-math}
\setmainfont{Latin Modern Roman}
\setmathfont{Latin Modern Math}
\usepackage{amsmath,amssymb,mathtools}
\usepackage{xurl,hyperref}
\hypersetup{colorlinks=true,urlcolor=blue,linkcolor=blue}
\setcounter{secnumdepth}{0}
\setlength{\parindent}{0pt}
\setlength{\parskip}{5pt}
\setlength{\emergencystretch}{3em}
\newcommand{\sref}[2]{\hyperlink{source:#1:#2}{[S#2]}}
\newcommand{\cataloguescope}[1]{\paragraph{Source question.}#1}
\newcommand{\floor}[1]{\left\lfloor#1\right\rfloor}
\begin{document}
% UnsolvedMath ID 30003377; source catalogue code OWR-15210-003
\setcounter{section}{30003376}
\section{Assumptions and Function Classes for Function Secret Sharing}\label{30003377}
{\small\sffamily UnsolvedMath ID 30003377 · Computer Science}
\subsection{Definitions and mathematical statement}

\paragraph{Shared notation.}Standard number systems are denoted by $\mathbb N,\mathbb Z,\mathbb Q,\mathbb R,\mathbb C$; index ranges, computational models and specialized notation follow the cited problem source and the definitions below.


Function secret sharing distributes keys whose local evaluations combine to the output of a function while hiding that function from the prescribed coalitions. A point function is nonzero at one designated input and zero elsewhere.

A cryptographic security requirement compares an efficient adversary\textquotesingle s advantage with a negligible function of the security parameter. A negligible function is smaller than every inverse polynomial eventually; computational and information-theoretic security are different requirements.
\paragraph{Statement recorded in the supplied source.}
Characterize the function classes for which efficient function secret-sharing schemes follow from one-way functions, and determine the necessary and sufficient cryptographic assumptions for useful FSS instances; in particular, improve one-way-function-based FSS for point functions with at least three parties. \sref{30003377}{2}
\paragraph{Scope qualification.}
The coalition threshold, leakage, output group, and efficiency parameters are source-specific.
\paragraph{Source status.}
The supplied catalogue labels this record partially solved. Its source-status metadata was checked against the supplied snapshot on 2026-10-11. This import does not independently establish current open status or mathematical correctness. The catalogue literature assessment is dated 2026-08-21.
\subsection{Short English statement}
Characterize which functions admit efficient secret sharing from one-way functions and improve the stated multiparty point-function case.
\subsection{Sources}
\begin{itemize}
\item[\hypertarget{source:30003377:1}{S1}] Ulam AI and the UnsolvedMath Contributors, UnsolvedMath: A Curated Collection of Open Mathematics Problems, record 30003377 (OWR-15210-003). Catalogue curation and metadata: CC BY 4.0. The supplied source question is retained; definitions and the English summary are editorial additions.
\url{https://huggingface.co/datasets/ulamai/UnsolvedMath}.
\item[\hypertarget{source:30003377:2}{S2}] Cryptography (2018), Oberwolfach Reports, DOI 10.4171/owr/2017/5, pp. 213--266
\url{https://doi.org/10.4171/owr/2017/5}.
\item[\hypertarget{source:30003377:3}{S3}] Boyle, Gilboa, and Ishai, Function Secret Sharing
\url{https://www.iacr.org/archive/eurocrypt2015/90560300/90560300.pdf}.
\item[\hypertarget{source:30003377:4}{S4}] Function Secret Sharing, Improvements and Extensions
\url{https://eprint.iacr.org/2018/707.pdf}.
\textit{Underlying Oberwolfach Reports material retains its source-specific licence, including CC BY-SA 4.0 where applicable; see THIRD\_PARTY\_NOTICES.md.}
\end{itemize}
\paragraph{Preparation and review.}
Original question attributed to the cited source and UnsolvedMath contributors. Codex (GPT-6; exact serving revision not exposed) assisted with definitions, English summaries, source inspection and formatting. No new human mathematical authorship or independent proof verification is asserted. Canonical ID 30003377 and source code OWR-15210-003 checked against the complete supplied registry and existing repository statements; known flagged repetitions and obvious identical conjectures excluded. Full mathematical-equivalence/literature review has not been performed. Statement preparation and final review followed contribution prompts C and D; automated checks verify structure and identity, not open status.
\label{end:30003377}
\end{document}
